% Chapter introduction: why the topic arose, who first asked, what it answers.
% STATUS: written 2026-11-05 (sources of the chapter read in full, each from the page)
Jacobson began the study of derivations as a piece of Lie theory in 1937.
His paper \emph{Abstract derivation and Lie algebras} \cite{Jacobson1937}
showed that the derivations of any algebra form a Lie algebra under the
commutator bracket.
It also asked how $\operatorname{Der} L$ reflects
the shape of $L$. The first answers came at the two extremes. A semisimple
Lie algebra in characteristic zero has only inner derivations
\cite[Theorem in \S 5.3, p.~23]{Humphreys1972}; a nilpotent Lie algebra
always has an outer one \cite[Exercise 10, p.~14]{Humphreys1972}. The
theory since then has walked between these poles, and the chapter follows
that walk.

Section~\ref{sec:lie-algebras:jacobson} reports the two extremal
answers.  Section~\ref{sec:lie-algebras:characteristically-nilpotent}
moves one step toward the nilpotent side: the characteristically nilpotent
Lie algebras of Dixmier and Lister, where every derivation is nilpotent.
On the other side, modular simple Lie algebras of Cartan and of
Albert--Zassenhaus type grow outer derivations that no characteristic-zero
analogue has; Section~\ref{sec:lie-algebras:simple} recounts this with
the classification theorem of Wilson and the computation of Winter of
$\operatorname{Der} L/\operatorname{ad} L$.  Section~\ref{sec:lie-algebras:vector-fields}
turns to the infinite Lie algebras of vector fields, where derivations
become concrete operators on polynomial functions.
Section~\ref{sec:lie-algebras:infinite-dimensional} closes with affine
Kac--Moody algebras, where the loop construction builds derivations into
the algebra itself. It also records the 2-local derivations of a reductive
Lie algebra, which fail to be derivations only when the centre is large
\cite[Corollary~1.3, p.~2]{ChenLuoXu2026}.

\hbfigure{14-lie-algebras/sl2-weights}{The adjoint action on $\mathfrak{sl}_2$. In characteristic $0$ every derivation of a semisimple Lie algebra is inner.}{fig:lie:sl2}
